\documentclass[11pt]{article}

\usepackage[margin=1in]{geometry}
\usepackage{booktabs}
\usepackage{tabularx}
\usepackage{enumitem}
\usepackage{listings}
\usepackage{xcolor}
\usepackage{hyperref}
\usepackage{fancyhdr}
\usepackage{array}
\usepackage{longtable}

\definecolor{codegray}{gray}{0.95}
\definecolor{commentgreen}{RGB}{0,110,60}
\definecolor{keywordblue}{RGB}{30,70,160}

\lstdefinestyle{pythonstyle}{
  language=Python,
  backgroundcolor=\color{codegray},
  basicstyle=\ttfamily\small,
  keywordstyle=\color{keywordblue}\bfseries,
  commentstyle=\color{commentgreen},
  stringstyle=\color{purple},
  showstringspaces=false,
  columns=fullflexible,
  keepspaces=true,
  frame=single,
  framerule=0.3pt,
  breaklines=true,
  tabsize=4
}
\lstset{style=pythonstyle}

\pagestyle{fancy}
\fancyhf{}
\lhead{Computational Thinking}
\rhead{Python: Expressions and Simple Scripts}
\cfoot{\thepage}

\hypersetup{
  colorlinks=true,
  linkcolor=blue,
  urlcolor=blue
}

\title{\textbf{Computational Thinking
%\\ Run Programs, Use an IDE or Notebook, Work with Expressions, and Write Simple Scripts
}}
\author{Alexander Abuabara}
\date{}

\begin{document}
\maketitle

\section*{Pedagogy}

This session introduces Python as a notation for computational problem solving. Students should leave with a basic mental model of how a computer evaluates expressions, stores values, executes code, and turns a small problem specification into a repeatable program.

This session reinforce the idea:

\begin{quote}
\textbf{Python is the notation; computation is the subject.}
\end{quote}

For each new piece of Python, ask:

\begin{enumerate}
    \item What problem are we representing?
    \item What values does the computer need?
    \item What operation should it perform?
    \item How can we tell whether the result makes sense?
\end{enumerate}

Therefore, this session establishes a clean conceptual foundation for later work on control flow.
It does not introduce conditionals, loops, functions, or collections unless time permits.

\section*{Learning Objectives}

By the end of the session, students should be able to:

\begin{enumerate}[label=\arabic*.]
    \item Explain, at a basic level, what it means for a computer to execute a program.
    \item Distinguish among Python code, an expression, a value, and a variable.
    \item Run Python interactively in a notebook or IDE.
    \item Predict and evaluate simple arithmetic expressions.
    \item Identify basic numerical types using \texttt{type()}.
    \item Assign values to variables and use variables in subsequent computations.
    \item Write, save, modify, and rerun a short Python program.
    \item Apply the computational-thinking cycle:
\end{enumerate}

\begin{center}
\textbf{problem $\rightarrow$ representation $\rightarrow$ instructions $\rightarrow$ execute $\rightarrow$ inspect $\rightarrow$ revise}
\end{center}

\section*{Preparation}

\begin{itemize}
    \item Access Python 3 in Jupyter, VS Code, Spyder, Positron, or another environment.
    \item Check the Jupyter notebook:
    \texttt{1\_computational\_thinking.ipynb}.
    \item If using notebooks, check how to run a cell.
    \item If using an IDE, check how to run both a single expression interactively.
\end{itemize}

\section*{Agenda}

\begin{enumerate}
\item What does it mean to compute?
\item Python as an executable language; notebook/IDE/script demonstration.
\item Predict, then run arithmetic expressions.
\item Objects, values, and numerical types.
\item Variables, assignment, and notebook state.
\item Build a tiny computational solution.
\item Debugging exercise.
\item Wrap-up and exit ticket.
\end{enumerate}

\section{What Does It Mean to Compute?}

\begin{quote}
A dataset contains 1,250 records. Processing one record takes approximately 0.018 seconds. Approximately how long will the entire dataset take to process?
\end{quote}

One minute to solve the problem manually.

\begin{itemize}
    \item What information did you use?
    \item What operation did you perform?
    \item What sequence of steps would you give a computer?
    \item What would change if the dataset contained 2.5 million records?
\end{itemize}

The instructional purpose is to distinguish solving a single instance manually from specifying a process that can be executed repeatedly.

\subsection*{Conceptual Model}

\begin{center}
\begin{tabular}{c}
Problem \\
$\downarrow$ \\
Inputs / Representation \\
$\downarrow$ \\
Operations / Algorithm \\
$\downarrow$ \\
Program \\
$\downarrow$ \\
Execution \\
$\downarrow$ \\
Output \\
$\downarrow$ \\
Check / Revise
\end{tabular}
\end{center}

Computational thinking means describing problems and solutions precisely enough that some or all of the solution can be carried out computationally.

\section{Python as an Executable Language}

Open the class notebook or Python console and run:

\begin{lstlisting}
2 + 3
\end{lstlisting}

Then:

\begin{lstlisting}
10 * 4
\end{lstlisting}

Ask: \emph{What did Python do?}

Introduce three working terms:

\begin{description}[style=nextline]
    \item[Expression] Code that evaluates to a value.
    \item[Value] The result produced by evaluating an expression.
    \item[Interpreter] Software that reads and executes/evaluates Python code.
\end{description}

\subsection*{Notebook Versus Script}

Notebook example:

\begin{lstlisting}
distance = 150
time = 3
distance / time
\end{lstlisting}

Saved script example:

\begin{lstlisting}
distance = 150
time = 3
speed = distance / time

print(speed)
\end{lstlisting}

Ask:

\begin{quote}
Why do we typically need \texttt{print()} in a script, while the final expression in a notebook cell may display automatically?
\end{quote}

Emphasize the conceptual distinction:

\begin{center}
\begin{tabular}{p{0.42\textwidth} p{0.42\textwidth}}
\toprule
\textbf{Interactive environment} & \textbf{Script} \\
\midrule
write $\rightarrow$ execute $\rightarrow$ inspect $\rightarrow$ modify $\rightarrow$ execute &
write a sequence of instructions $\rightarrow$ execute the program \\
\bottomrule
\end{tabular}
\end{center}

\section{Predict, Then Run}

Students open the provided notebook or an interactive Python console.

\textbf{Rule for the activity:}

\begin{quote}
Do not run an expression until you have predicted its result.
\end{quote}

This turns the interpreter into a tool for checking reasoning rather than trial-and-error guessing.

\subsection*{Arithmetic Expressions}

Students predict and then execute:

\begin{lstlisting}
3 + 4
3 * 4
3 + 4 * 2
(3 + 4) * 2
10 / 4
10 // 4
10 % 4
2 ** 5
\end{lstlisting}

Discuss:

\begin{itemize}
    \item operator precedence,
    \item parentheses,
    \item \texttt{/} versus \texttt{//},
    \item remainder with \texttt{\%},
    \item exponentiation with \texttt{**}.
\end{itemize}

Ask students to translate expressions into words. For example, \texttt{10 \% 4} can be interpreted as finding the remainder after dividing 10 into groups of 4.

\section{Objects, Values, and Types}

Run:

\begin{lstlisting}
type(3)
type(3.0)
type(10 / 2)
\end{lstlisting}

Ask:

\begin{quote}
Are \texttt{3} and \texttt{3.0} the same thing computationally?
\end{quote}

They may represent the same mathematical quantity, but Python stores and treats them as different numerical types.

\subsection*{Working Vocabulary}

\begin{center}
\begin{tabular}{ll}
\toprule
\texttt{3} & integer value (\texttt{int}) \\
\texttt{3.0} & floating-point value (\texttt{float}) \\
\texttt{2 + 3} & expression \\
\texttt{5} & value produced by the expression \\
\bottomrule
\end{tabular}
\end{center}

Then ask students to predict both value and type:

\begin{lstlisting}
2 + 3.0
type(2 + 3.0)
\end{lstlisting}

\subsection*{Quick Checkpoint}

What will Python report?

\begin{lstlisting}
type(7 / 2)
\end{lstlisting}

\begin{enumerate}[label=\Alph*.]
    \item \texttt{int}
    \item \texttt{float}
    \item error
    \item impossible to determine
\end{enumerate}

Students vote before the code is executed.

\section{Variables and Assignment}

Introduce:

\begin{lstlisting}
records = 1250
seconds_per_record = 0.018
\end{lstlisting}

Then:

\begin{lstlisting}
total_seconds = records * seconds_per_record
total_seconds
\end{lstlisting}

Avoid initially describing \texttt{=} as mathematical equality. Instead:

\begin{quote}
\texttt{records = 1250} associates the name \texttt{records} with the value \texttt{1250}.
\end{quote}

Demonstrate reassignment:

\begin{lstlisting}
records = 1250
records

records = 2500
records
\end{lstlisting}

\subsection*{Notebook State}

Demonstrate:

\begin{lstlisting}
x = 10
\end{lstlisting}

Then, in a separate cell:

\begin{lstlisting}
x * 2
\end{lstlisting}

Edit the first cell to:

\begin{lstlisting}
x = 100
\end{lstlisting}

but do not run it. Run the second cell again and discuss the result.

\textbf{Key message:}

\begin{quote}
Computers execute what you run, not what you intended to run.
\end{quote}

Use this moment to explain that notebook state depends on execution history, not merely visible cell order.

\section{Build a Tiny Computational Solution}

\subsection*{Lab Problem: Processing-Time Estimator}

\textbf{Scenario:} A data-processing job has a known number of records and an estimated processing time per record. Write a short Python program that estimates the total processing time.

\subsubsection*{Step 1: Identify Inputs and Output}

Students write:

\begin{verbatim}
Input:
    number of records
    seconds per record

Output:
    estimated processing time
\end{verbatim}

Then:

\begin{lstlisting}
records = 1250
seconds_per_record = 0.018
\end{lstlisting}

\subsubsection*{Step 2: Compute}

\begin{lstlisting}
total_seconds = records * seconds_per_record
\end{lstlisting}

\subsubsection*{Step 3: Display}

\begin{lstlisting}
print(total_seconds)
\end{lstlisting}

\subsubsection*{Step 4: Improve the Representation}

\begin{lstlisting}
total_minutes = total_seconds / 60
print(total_minutes)
\end{lstlisting}

\subsubsection*{Step 5: Test with Another Case}

Change:

\begin{lstlisting}
records = 1_000_000
\end{lstlisting}

Rerun the program.

Ask:

\begin{quote}
Which parts of your program represent the particular problem instance, and which parts represent the general computational solution?
\end{quote}

Desired distinction:

\begin{itemize}
    \item \textbf{Data/parameters:} \texttt{records}, \texttt{seconds\_per\_record}
    \item \textbf{Computation:} \texttt{records * seconds\_per\_record}
\end{itemize}

\subsection*{Lab Challenge A: Predict Before Executing}

Before running:

\begin{lstlisting}
records = 100
seconds_per_record = 0.25

total_seconds = records * seconds_per_record
total_minutes = total_seconds / 60
\end{lstlisting}

Students record predictions for \texttt{total\_seconds} and \texttt{total\_minutes}, then verify them.

\subsection*{Lab Challenge B: Modify the Model}

Add:

\begin{lstlisting}
processors = 4
parallel_seconds = total_seconds / processors
\end{lstlisting}

Ask:

\begin{quote}
What assumption is the program making?
\end{quote}

Important answer: it assumes ideal scaling, so four processors provide a fourfold speedup.

\textbf{Computational-thinking lesson:} A program implements a model, and the model contains assumptions.

\subsection*{Lab Challenge C: Write a Small Script}

Students choose one:

\begin{itemize}
    \item Fahrenheit $\rightarrow$ Celsius
    \item miles $\rightarrow$ kilometers
    \item number of files $\times$ average file size $\rightarrow$ total storage
    \item records / elapsed seconds $\rightarrow$ records per second
\end{itemize}

Requirements:

\begin{itemize}
    \item at least two variables,
    \item at least two expressions,
    \item at least one newly computed variable,
    \item at least one \texttt{print()} statement,
    \item meaningful variable names.
\end{itemize}

Avoid:

\begin{lstlisting}
x = 100
y = 5
z = x / y
\end{lstlisting}

Prefer:

\begin{lstlisting}
records_processed = 100
elapsed_seconds = 5
records_per_second = records_processed / elapsed_seconds

print(records_per_second)
\end{lstlisting}

\section{Debugging Micro-Exercise \hfill 65--70 min}

Give students these three examples.

\subsection*{Example A}

\begin{lstlisting}
records = 1000
seconds = 20
rate = records / Seconds
print(rate)
\end{lstlisting}

\subsection*{Example B}

\begin{lstlisting}
records = 1000
seconds = 20
rate = records seconds
print(rate)
\end{lstlisting}

\subsection*{Example C}

\begin{lstlisting}
records = 1000
seconds = 20
rate = seconds / records
print(rate)
\end{lstlisting}

Ask students to classify the issue:

\begin{itemize}
    \item A: name/case error,
    \item B: syntax error,
    \item C: program runs, but the computation is logically wrong.
\end{itemize}

Emphasize that successful execution does not guarantee a correct program.

Establish the recurring course habit:

\begin{center}
\textbf{Predict $\rightarrow$ Run $\rightarrow$ Inspect $\rightarrow$ Explain}
\end{center}

rather than:

\begin{center}
Run $\rightarrow$ hope
\end{center}

\section{Wrap-Up and Exit Ticket \hfill 70--75 min}

Return to the opening problem and ask students to describe the solution without Python:

\begin{enumerate}
    \item Represent the number of records.
    \item Represent the processing time per record.
    \item Multiply them.
    \item Report the result.
\end{enumerate}

Then map those ideas to Python:

\begin{lstlisting}
records = 1250
seconds_per_record = 0.018

total_seconds = records * seconds_per_record

print(total_seconds)
\end{lstlisting}

\subsection*{Exit Ticket}

\textbf{1. Predict}

\begin{lstlisting}
x = 8
y = 3
z = x / y
\end{lstlisting}

What is the type of \texttt{z}?

\medskip

\textbf{2. Explain}

In one sentence, explain the difference between:

\begin{lstlisting}
3 + 4
\end{lstlisting}

and:

\begin{lstlisting}
result = 3 + 4
\end{lstlisting}

\medskip

\textbf{3. Computational Thinking}

For your lab program, identify:

\begin{verbatim}
Input(s):
Operation(s):
Output:
One assumption:
\end{verbatim}

\section*{Suggested Instructor Notebook Structure}

\begin{verbatim}
00 - Today's question
     "How do we turn a simple problem into something
      a computer can execute?"

01 - First Python expressions
02 - Predict -> run -> explain
03 - Numerical types
04 - Variables and assignment
05 - Notebook state
06 - Guided lab: processing-time estimator
07 - Lab challenges
08 - Debugging exercise
09 - Exit ticket
\end{verbatim}

The notebook should remain activity-oriented rather than becoming a page of lecture notes. Each conceptual section should be followed quickly by executable work.

\section*{Core Vocabulary}

\begin{longtable}{p{1.4in} p{4.9in}}
\toprule
\textbf{Term} & \textbf{Working Definition} \\
\midrule
Program & A sequence of instructions that a computer can execute. \\
Expression & Python code that evaluates to a value. \\
Value & A piece of data produced or manipulated by a program. \\
Type & A category that determines how a value is represented and can be used. \\
Variable & A name associated with a value. \\
Assignment & Associating a name with a value. \\
Interpreter & Software that executes or evaluates Python instructions. \\
Script & Python code saved as a program and executed. \\
Notebook & An interactive document containing executable code cells and explanatory content. \\
\bottomrule
\end{longtable}

\section*{Extra}

Convert this word problem into a program:

\begin{quote}
A sensor generates 240 observations per minute. Each observation requires 16 bytes. Estimate the storage required for one hour of observations.
\end{quote}

Before coding, require:

\begin{verbatim}
Inputs:
Operations:
Output:
\end{verbatim}

Then have students code, test, and exchange programs with another pair.
The receiving pair changes one input and predicts the new output before running the program.

This provides a natural transition to the next class question:

\begin{quote}
Once we can represent data and evaluate expressions, how can a program make decisions or repeat work?
\end{quote}

\begin{thebibliography}{9}

\bibitem{guttag}
John V. Guttag.
\newblock \emph{Introduction to Computation and Programming Using Python: With Application to Computational Modeling and Understanding Data}.
\newblock Third edition, MIT Press.
\newblock \url{https://mitpress.mit.edu/9780262542364/introduction-to-computation-and-programming-using-python/}

\end{thebibliography}

\end{document}
